\section*{Chapitre \ref{ch:hf:etf-market-making} --- Market making sur ETF}

\begin{solution}{exo:hf:etf-market-making:1}
$1\,250/(50\,000\times291.58)=0.86$ points de base.
\end{solution}

\begin{solution}{exo:hf:etf-market-making:2}
$100\times1.015\times0.995=100.99$.
\end{solution}

\begin{solution}{exo:hf:etf-market-making:3}
L'écart vaut $0.05\times50\,000=\$2\,500$~; le coût vaut $\$1\,250+0.0002\times100\times50\,000=\$2\,250$. Il faut créer, avec un edge de 250~\$ par unité.
\end{solution}

\begin{solution}{exo:hf:etf-market-making:4}
La valeur liquidative se situait 557 points de base au-dessus de la valeur réelle des obligations (obsolescence des prix)~; le mid se situait 5 points de base
en dessous de la valeur réelle (la décote du teneur de marché). Seule la seconde est négociable contre les obligations, et seulement par qui peut vendre les
obligations à leur valeur réelle, ce que, dans un gel, personne ne peut faire à grande taille.
\end{solution}

\begin{solution}{exo:hf:etf-market-making:5}
Sa cotation s'écarte en moyenne de 13.7 points de base de la valeur réelle, soit davantage que sa demi-fourchette~: les clients qui voient le proxy achètent
lorsque son cours vendeur est inférieur à la valeur réelle et vendent lorsque son cours acheteur est supérieur, et le flux qu'il reçoit est celui que ses erreurs attirent. Il gagne 132\,100~\$ un jour normal, contre 163\,600~\$.
\end{solution}

\begin{solution}{exo:hf:etf-market-making:6}
Le rachat livre des obligations qui doivent être vendues à la fourchette du gel, six fois la normale~: 262\,000~\$ sur cinq jours, pour un inventaire
déjà couvert dans le proxy et qui aurait été dénoué par les cotations après le gel.
\end{solution}

\begin{solution}{exo:hf:etf-market-making:7}
L'obsolescence maximale de la valeur liquidative passe de 557 à 722 points de base, l'erreur du \omterm{def:hf:etf-market-making:pricing}{panier de valorisation} pendant le gel de 13.8 à 16.5 points de base~;
le teneur de marché sur panier gagne encore 3.48 millions de dollars (contre 3.43 millions) et le teneur de marché sur valeur liquidative perd 96.9 millions de dollars (contre 84.7 millions).
L'obsolescence nuit à celui qui cote à partir de prix périmés~; un panier piloté par le proxy se dégrade lentement.
\end{solution}

\begin{solution}{exo:hf:etf-market-making:8}
La valeur liquidative était obsolète~: les obligations valaient moins qu'elle ne l'indiquait. Le rachat livre les obligations, qui doivent être vendues
à leur valeur réelle et à la fourchette de stress~; l'essentiel des 5~\% était l'erreur de la valeur liquidative, et non un profit.
L'arbitrage porte sur la valeur exécutable des obligations, et non sur le chiffre publié.
\end{solution}

\begin{solution}{pb:hf:etf-market-making:1}
\begin{enumerate}
\item Voir la \cref{def:hf:etf-market-making:pricing}~; $\hat P_{i,t}=P_i(1+\beta_i(G_t-G_{t_i}))$.
\item Chaque position est obsolète~: la clôture date d'un certain temps, auquel s'ajoutent le mouvement du contrat à terme overnight et le mouvement de la devise.
\item Les écarts sont généralement de l'ordre de 200 points de base, davantage pour les fonds internationaux et peu liquides,
encore environ 100 après contrôle de l'obsolescence des prix.
\item Valeur liquidative 13.7 et 214 points de base~; \omterm{def:hf:etf-market-making:pricing}{panier de valorisation} 3.5 et 9.9~; dernières cotations âgées en moyenne de 2 et 11 heures de séance.
\item Voir la \cref{def:hf:etf-market-making:fee}~; 1\,250~\$ par ordre sur une unité de 14.6 millions de dollars, jusqu'à 3.0~\% de plus pour les créations en espèces.
\item Voir la \cref{def:hf:etf-market-making:custom}~; Rule 6c-11, avec des politiques et procédures écrites.
\item Créer si $(P^{\mathrm{ETF}}-V)N>F+c_{\mathrm{buy}}VN$~; racheter si $(V-P^{\mathrm{ETF}})N>F+c_{\mathrm{sell}}VN$.
\item Les obligations reçues doivent être vendues à la fourchette de stress, et le fonds peut choisir les obligations à livrer.
\item Teneur de marché sur valeur liquidative 111\,600~\$ (écart type 16\,000~\$)~; teneur de marché sur panier 149\,100~\$ (écart type 9\,300~\$).
\item $-\$69.9$ millions et $+\$3.55$ millions~; 95~\% de la perte sur les transactions avec des clients à des prix supérieurs à la valeur des obligations.
\item Environ 14\,500~\$ un jour normal~; un tiers de moins que l'écart type du jour.
\item 262\,000~\$ en fourchettes d'obligations.
\item Le teneur de marché sur panier a gagné 3.2 à 3.6 millions de dollars pendant le gel, tout en cotant jusqu'à 532 points de base en dessous
de la valeur liquidative, dont 557 relevaient de l'obsolescence et $-5$ de sa propre décote par rapport à la valeur réelle.
\item Le teneur de marché qui dispose d'un \omterm{def:hf:etf-market-making:pricing}{panier de valorisation} et d'une couverture par proxy, à environ 5 à 10 points de base en dessous de la valeur réelle.
\item Un teneur de marché qui doit cesser d'acheter laisse les actions baisser sous la valeur réelle jusqu'à l'arrivée d'autres acheteurs~:
la décote devient un prix pour le bilan.
\item Lorsque les actions se négocient sous la valeur à laquelle le panier peut être vendu, après frais et fourchette, et que le rachat peut être mené à son terme.
\item L'âge des prix qui les sous-tendent, et une valeur ajustée à partir des proxys observables.
\item Régresser les rendements des obligations entre deux cotations sur le rendement du proxy sur le même intervalle, regroupés par notation et échéance.
\item La cotation de taux sur une valeur périmée~: elle se situe là où l'obsolescence est la plus forte.
\item Son propre \omterm{def:hf:etf-market-making:pricing}{panier de valorisation}, adossé à sa couverture et à son bilan.
\end{enumerate}
\end{solution}

\begin{solution}{iq:hf:etf-market-making:1}
Acheter le fonds et vendre le panier, ou acheter et racheter~; tout le monde le fait aussi, si bien que l'écart se referme avant qu'une unité puisse être
constituée, sauf s'il tient dans les coûts des fourchettes et des frais du panier.
\iqlookfor{valeur exécutable du panier et concurrence.}
\end{solution}

\begin{solution}{iq:hf:etf-market-making:2}
Déplacer la dernière transaction de chaque obligation selon sa sensibilité aux proxys liquides (contrats à terme sur emprunts d'État, indices de crédit, fonds similaires)
depuis cette transaction~; utiliser les cotations de dealers et les données des plateformes lorsqu'elles existent~; pondérer selon les positions du fonds~; vérifier par comparaison avec des fonds similaires.
\iqlookfor{un \omterm{def:hf:etf-market-making:pricing}{panier de valorisation} ajusté par proxys.}
\end{solution}

\begin{solution}{iq:hf:etf-market-making:3}
Les fourchettes exécutables des obligations, les conditions de rachat et le panier du fonds, la liquidité et la base du proxy, le bilan et ses limites~;
en général, on couvre dans le proxy et on conserve la position, en rachetant lorsque les obligations peuvent être vendues.
\iqlookfor{coûts de chaque sortie et de la couverture.}
\end{solution}

\begin{solution}{iq:hf:etf-market-making:4}
En situation de stress, la valeur liquidative est obsolète~: la valoriser à ce cours montre le livre à une valeur que personne ne peut réaliser, masque la perte
et incite à acheter davantage. Valoriser à une valeur de modèle issue du \omterm{def:hf:etf-market-making:pricing}{panier de valorisation}, et signaler l'écart.
\iqlookfor{valorisations indépendantes.}
\end{solution}

\begin{solution}{iq:hf:etf-market-making:5}
Fichiers de détention chargés avant l'ouverture~; un cache partagé des cours des instruments et des mouvements des proxys, mis à jour à chaque tick~; la valeur
de chaque fonds mise à jour de façon incrémentale à partir des instruments qui ont bougé~; des modèles de prix obsolètes exécutés sur un minuteur~; publication vers les moteurs de cotation avec horodatage et âge.
\iqlookfor{calcul incrémental et état partagé.}
\end{solution}

\begin{solution}{iq:hf:etf-market-making:6}
Si le rendement de la position depuis sa dernière cotation vaut $\beta\Delta G+\varepsilon$ avec $\varepsilon$ non corrélé à $\Delta G$, la prévision des
moindres carrés sachant $\Delta G$ est $\beta\Delta G$. Elle est biaisée lorsque $\beta$ est mal estimé, lorsque le proxy a sa propre base, ou lorsque les
nouvelles concernant la position sont corrélées à l'absence de transactions.
\iqlookfor{l'argument de la régression et ses modes de défaillance.}
\end{solution}
